% !TeX program = xelatex
% UTF-8; compile twice with XeLaTeX.
% Source: 问题一_定位区域模型.md
% This file preserves the supplied text, mathematics, data and section numbering.
\documentclass[UTF8,a4paper,zihao=-4,fontset=fandol]{ctexart}
\XeTeXgenerateactualtext=1
\usepackage[left=23mm,right=23mm,top=22mm,bottom=23mm,headheight=15pt,headsep=7mm]{geometry}
\usepackage{amsmath,amssymb,mathtools}
\usepackage{booktabs,tabularx,array}
\usepackage{enumitem}
\usepackage{needspace}
\usepackage{tcolorbox}
\usepackage{fancyhdr}
\usepackage[unicode,hidelinks]{hyperref}
\usepackage{bookmark}
\hypersetup{pdftitle={基于角域交会与最小包围圆的无源测向快速定位模型}}
\renewcommand{\thesection}{5.\arabic{section}}
\renewcommand{\thesubsection}{\thesection.\arabic{subsection}}
\ctexset{
  section={format=\large\bfseries,beforeskip=1.2em,afterskip=0.55em},
  subsection={format=\normalsize\bfseries,beforeskip=0.9em,afterskip=0.35em}
}
\linespread{1.16}
\setlength{\parskip}{0.25em}
\setlength{\emergencystretch}{2em}
\setlist[enumerate]{label=\arabic*.,leftmargin=2em,itemsep=0.15em,topsep=0.4em,parsep=0pt}
\setlist[itemize]{leftmargin=2em,itemsep=0.15em,topsep=0.4em,parsep=0pt}
\pagestyle{fancy}
\fancyhf{}
\fancyhead[L]{\small 无源测向快速定位模型}
\fancyfoot[C]{\small\thepage}
\renewcommand{\headrulewidth}{0.3pt}
\renewcommand{\footrulewidth}{0pt}
\fancypagestyle{plain}{\fancyhf{}\fancyfoot[C]{\small\thepage}\renewcommand{\headrulewidth}{0pt}}
\widowpenalty=10000
\clubpenalty=10000
\displaywidowpenalty=10000
\raggedbottom

% BEGIN PROCESS DRAWING DEFINITIONS
\def\ThirdIntersectionPath{(2.5925925926,1.6888888889)--(1.2380952381,1.4857142857)--(1.0357142857,1.2428571429)--(2.8437500000,1.3875000000)--cycle}
% Process illustrations: geometric steps, not numerical result charts.
\usepackage{tikz}
\usepackage{caption}
\usetikzlibrary{arrows.meta,calc,patterns}
\definecolor{procblue}{HTML}{3976A3}
\definecolor{procteal}{HTML}{187D79}
\definecolor{procorange}{HTML}{B77A36}
\definecolor{procred}{HTML}{BA4D43}
\captionsetup{font=small,labelfont=bf,labelsep=quad,skip=7pt,hypcap=false}
\renewcommand{\figurename}{图}
\tikzset{
 proc/.style={>=Stealth,line cap=round,line join=round,font=\footnotesize},
 ptitle/.style={font=\small\bfseries,align=center},
 pnote/.style={font=\footnotesize,align=center,text width=4.5cm},
 edge/.style={procteal,thick},
 kept/.style={procteal,fill=procteal},
 rejected/.style={procred,thick}
}
\newcommand{\ProcArrow}[1]{\draw[->,black!45,thick] (#1,2.0)--++(0.5,0);}
\newcommand{\ProcCross}[2]{\draw[rejected] (#1-0.09,#2-0.09)--(#1+0.09,#2+0.09);\draw[rejected] (#1-0.09,#2+0.09)--(#1+0.09,#2-0.09);}
\newcommand{\ProcPentagon}{(0,0)--(4,0)--(4,1)--(2,3)--(0,3)--cycle}
\newcommand{\ProcBounds}{%
 \draw[black!30] (0,-.35)--(0,3.75) (4,-.35)--(4,3.75);
 \draw[black!30] (-.3,0)--(5.4,0) (-.3,3)--(5.4,3);
 \draw[black!30] (1.25,3.75)--(5.35,-.35) (3.65,-.35)--(5.35,1.35);
}

\newcommand{\DirectionToConstraints}{%
\begin{tikzpicture}[proc]
 \path[use as bounding box] (-.15,-.55) rectangle (15.05,4.0);
 \node[ptitle] at (2.15,3.8) {(a) 读到一个方向};
 \node[ptitle] at (7.45,3.8) {(b) 展开误差范围};
 \node[ptitle] at (12.75,3.8) {(c) 写成两个半平面};
 \begin{scope}[shift={(.2,.55)}]
  \draw[->,black!40] (0,.3)--(3.85,.3);
  \draw[->,procblue,very thick] (0,.3)--(3.65,1.774495);
  \fill (0,.3) circle (1.7pt) node[below] {$S_i$};
  \draw[->,procblue] (.85,.3) arc (0:22:.85);
  \node[procblue] at (1.18,.53) {$\theta_i$};
  \node[procblue,above] at (2.35,1.4) {示向度};
 \end{scope}
 \begin{scope}[shift={(5.5,.55)}]
  \fill[procteal!12] (0,.3)--(3.75,.9)--(3.2,2.55)--cycle;
  \draw[->,procblue,thick] (0,.3)--(3.75,.9);
  \draw[->,procblue,thick] (0,.3)--(3.2,2.55);
  \draw[->,procblue,dashed] (0,.3)--(3.65,1.774495);
  \fill (0,.3) circle (1.7pt) node[below] {$S_i$};
  \draw[<->,procorange] (9:1.85)++(0,.3) arc (9:22:1.85);
  \draw[<->,procorange] (22:1.85)++(0,.3) arc (22:35:1.85);
  \node[procorange] at (2.17,.8) {$\varepsilon$};
  \node[procorange] at (1.97,1.63) {$\varepsilon$};
  \node[procteal] at (2.85,1.45) {$W_i$};
 \end{scope}
 \begin{scope}[shift={(10.8,.55)}]
  \begin{scope}
   \clip (-.1,.05) rectangle (3.95,2.65);
   \fill[procblue!10] (-.1,.284)--(3.95,.932)--(3.95,2.65)--(-.1,2.65)--cycle;
   \fill[procorange!11] (-.1,.23)--(3.95,3.065)--(3.95,.05)--(-.1,.05)--cycle;
   \fill[procteal!25] (0,.3)--(3.95,.932)--(3.95,3.065)--cycle;
   \draw[procblue,thick] (-.1,.284)--(3.95,.932);
   \draw[procorange,thick] (-.1,.23)--(3.95,3.065);
  \end{scope}
  \fill (0,.3) circle (1.7pt) node[below] {$S_i$};
  \node[procblue] at (1.15,2.35) {$H_i^-$};
  \node[procorange] at (2.9,.28) {$H_i^+$};
  \node[procteal,font=\scriptsize] at (2.15,1.5) {$W_i=H_i^-\cap H_i^+$};
 \end{scope}
 \ProcArrow{4.55}\ProcArrow{9.85}
 \node[pnote] at (2.15,-.13) {方向给出了方位，\textbf{尚未给出距离}};
 \node[pnote] at (7.45,-.13) {向两侧各展开 $\varepsilon$，\textbf{得到候选角域}};
 \node[pnote] at (12.75,-.13) {两条不等式同时成立，\textbf{才保留该点}};
\end{tikzpicture}%
}

\newcommand{\ProgressiveIntersection}{%
\begin{tikzpicture}[proc]
 \path[use as bounding box] (-.15,-.55) rectangle (15.05,4.0);
 \node[ptitle] at (2.15,3.8) {(a) 第一次：一个角域};
 \node[ptitle] at (7.45,3.8) {(b) 第二次：保留交集};
 \node[ptitle] at (12.75,3.8) {(c) 再测一次：继续收缩};
 \foreach \off in {0,5.3,10.6}{
  \begin{scope}[shift={(\off+.15,.55)}]
   \draw[black!12] (0,0) rectangle (4,2.7);
  \end{scope}
 }
 \begin{scope}[shift={(.15,.55)}]
  \begin{scope}\clip (0,0) rectangle (4,2.7);
   \fill[procblue!13] (0,0)--(4,1.2)--(4,4.8)--cycle;
   \draw[procblue,thick] (0,0)--(4,1.2) (0,0)--(4,4.8);
  \end{scope}
  \fill (0,0) circle (1.5pt) node[below left] {$S_1$};
  \node[procblue] at (2.8,1.8) {$W_1$};
 \end{scope}
 \begin{scope}[shift={(5.45,.55)}]
  \begin{scope}\clip (0,0) rectangle (4,2.7);
   \fill[procblue!8] (0,0)--(4,1.2)--(4,4.8)--cycle;
   \fill[procorange!9] (4,0)--(0,1.2)--(0,4.8)--cycle;
   \draw[procblue!70] (0,0)--(4,1.2) (0,0)--(4,4.8);
   \draw[procorange!80] (4,0)--(0,1.2) (4,0)--(0,4.8);
   \filldraw[edge,fill=procteal!22] (2,.6)--(3.2,.96)--(2,2.4)--(.8,.96)--cycle;
  \end{scope}
  \fill (0,0) circle (1.5pt) node[below left] {$S_1$};
  \fill (4,0) circle (1.5pt) node[below right] {$S_2$};
  \node[procteal] at (2,1.4) {$P_2$};
 \end{scope}
 \begin{scope}[shift={(10.75,.55)}]
  \filldraw[black!25,fill=black!5] (2,.6)--(3.2,.96)--(2,2.4)--(.8,.96)--cycle;
  \fill[procblue!12] (0,1.16)--(4,1.48)--(4,1.9)--(0,1.3)--cycle;
  \draw[procblue] (0,1.16)--(4,1.48) (0,1.3)--(4,1.9);
  \filldraw[edge,fill=procteal!30] \ThirdIntersectionPath;
  \node[procblue,anchor=west] at (.05,1.55) {$W_3$};
  \node[procteal,fill=white,inner sep=1pt] at (2.15,1.44) {$P_3$};
  \node[black!50] at (2.05,2.18) {$P_2$};
 \end{scope}
 \ProcArrow{4.55}\ProcArrow{9.85}
 \node[pnote] at (2.15,-.13) {$P_1=W_1$\\只排除角域以外的位置};
 \node[pnote] at (7.45,-.13) {$P_2=P_1\cap W_2$\\必须同时符合两次测向};
 \node[pnote] at (12.75,-.13) {$P_3=P_2\cap W_3\subseteq P_2$\\新信息不能扩大候选区域};
\end{tikzpicture}%
}

\newcommand{\FeasibilityCases}{%
\begin{tikzpicture}[proc]
 \path[use as bounding box] (-.15,-.55) rectangle (15.05,3.75);
 \node[ptitle] at (2.15,3.55) {(a) 空集：约束矛盾};
 \node[ptitle] at (7.45,3.55) {(b) 无界：仍能走到任意远};
 \node[ptitle] at (12.75,3.55) {(c) 有界：形成定位多边形};
 \begin{scope}[shift={(.15,.6)}]
  \fill[procblue!13] (0,0) rectangle (1.3,2.2);
  \fill[procorange!14] (2.7,0) rectangle (4,2.2);
  \draw[procblue,thick] (1.3,0)--(1.3,2.2);
  \draw[procorange,thick] (2.7,0)--(2.7,2.2);
  \draw[procblue,->] (1.2,1.1)--(.35,1.1);
  \draw[procorange,->] (2.8,1.1)--(3.65,1.1);
  \node[procred] at (2,1.5) {$\varnothing$};
  \node at (2,.45) {没有公共部分};
 \end{scope}
 \begin{scope}[shift={(5.45,.6)}]
  \fill[procteal!12] (.25,.45)--(4,.15)--(4,2.25)--cycle;
  \draw[edge,->] (.25,.45)--(4,.15);
  \draw[edge,->] (.25,.45)--(4,2.25);
  \draw[procorange,very thick,->] (1.1,.75)--(3.8,1.12) node[above left] {可无限延伸};
  \fill (.25,.45) circle (1.5pt) node[below right] {有限顶点};
 \end{scope}
 \begin{scope}[shift={(10.95,.6)},x=.7cm,y=.7cm]
  \filldraw[edge,fill=procteal!14] \ProcPentagon;
  \draw[<->,black!50] (0,-.35)--(4,-.35);
  \draw[<->,black!50] (-.35,0)--(-.35,3);
  \node[procteal] at (1.8,1.3) {$P$};
 \end{scope}
 \node[pnote] at (2.15,-.13) {线性规划不可行\\$\Downarrow$ 报告数据异常};
 \node[pnote] at (7.45,-.13) {至少一个坐标方向无界\\$\Downarrow$ 补充测向};
 \node[pnote] at (12.75,-.13) {四个坐标方向的极值均有限\\$\Downarrow$ 继续求顶点与直径};
\end{tikzpicture}%
}

\newcommand{\VertexExtraction}{%
\begin{tikzpicture}[proc]
 \path[use as bounding box] (-.15,-.65) rectangle (15.05,4.1);
 \node[ptitle] at (2.15,3.85) {(a) 枚举候选交点};
 \node[ptitle] at (7.45,3.85) {(b) 检验全部约束};
 \node[ptitle] at (12.75,3.85) {(c) 只保留互异顶点};
 \begin{scope}[shift={(.25,.65)},x=.7cm,y=.7cm]
  \ProcBounds
  \foreach \p in {(0,0),(0,3),(4,0),(4,3),(4,1),(2,3),(5,0),(4.5,.5)}{
   \draw[procorange,fill=white] \p circle (1.8pt);
  }
 \end{scope}
 \begin{scope}[shift={(5.55,.65)},x=.7cm,y=.7cm]
  \ProcBounds
  \filldraw[edge,fill=procteal!12] \ProcPentagon;
  \foreach \p in {(0,0),(0,3),(4,0),(4,1),(2,3)}{\fill[procteal] \p circle (1.9pt);}
  \ProcCross{4}{3}\ProcCross{5}{0}\ProcCross{4.5}{.5}
  \node[procred,anchor=east,font=\scriptsize] at (5.3,3.5) {不满足约束};
 \end{scope}
 \begin{scope}[shift={(11.0,.65)},x=.7cm,y=.7cm]
  \filldraw[edge,fill=procteal!12] \ProcPentagon;
  \fill[procteal] (0,0) circle (1.9pt) node[below left] {$v_1$};
  \fill[procteal] (4,0) circle (1.9pt) node[below right] {$v_2$};
  \fill[procteal] (4,1) circle (1.9pt) node[right] {$v_3$};
  \fill[procteal] (2,3) circle (1.9pt) node[above] {$v_4$};
  \fill[procteal] (0,3) circle (1.9pt) node[above left] {$v_5$};
  \draw[black!45,->] (2.05,.8)--(3.85,.06);
  \node[align=center,font=\scriptsize] at (1.7,1.22) {重复交点\\只记一次};
 \end{scope}
 \ProcArrow{4.55}\ProcArrow{9.85}
 \node[pnote] at (2.15,-.23) {对非平行边界两两求交\\交点暂时都只是候选点};
 \node[pnote] at (7.45,-.23) {代回全部不等式\\有任一不满足，就剔除};
 \node[pnote] at (12.75,-.23) {合并重合点\\得到 $V=\{v_1,\ldots,v_h\}$};
\end{tikzpicture}%
}

\newcommand{\DiameterReduction}{%
\begin{tikzpicture}[proc]
 \path[use as bounding box] (-.15,-.45) rectangle (15.05,4.45);
 \node[ptitle] at (3.15,4.15) {(a) 定义：在整个区域中选两点};
 \node[ptitle] at (11.55,4.15) {(b) 计算：只比较所有顶点对};
 \begin{scope}[shift={(1.4,.65)},x=.9cm,y=.9cm]
  \filldraw[edge,fill=procteal!12] \ProcPentagon;
  \draw[black!50,dashed,thick] (.7,1.6)--(3.3,.6);
  \fill (.7,1.6) circle (1.6pt) node[above left] {$X$};
  \fill (3.3,.6) circle (1.6pt) node[above right] {$Y$};
  \draw[->,black!40] (.75,1.65) to[bend left] (1.45,2.2);
  \draw[->,black!40] (3.2,.65) to[bend right] (2.7,1.25);
 \end{scope}
 \draw[->,black!50,thick] (6.2,2.1)--(8.25,2.1) node[midway,above] {利用凸性};
 \begin{scope}[shift={(9.8,.65)},x=.9cm,y=.9cm]
  \fill[procteal!8] \ProcPentagon;
  \coordinate (a) at (0,0);\coordinate (b) at (4,0);\coordinate (c) at (4,1);
  \coordinate (d) at (2,3);\coordinate (e) at (0,3);
  \foreach \s/\t in {a/b,a/c,a/d,a/e,b/c,b/d,b/e,c/d,c/e,d/e}{\draw[black!25] (\s)--(\t);}
  \draw[edge] \ProcPentagon;
  \draw[procred,very thick] (b)--(e) node[midway,fill=white,inner sep=1pt] {$D$};
  \foreach \p in {a,b,c,d,e}{\fill[procteal] (\p) circle (1.7pt);}
  \node[procred,below right] at (b) {$A$};\node[procred,above left] at (e) {$B$};
 \end{scope}
 \node[pnote,text width=6.2cm] at (3.15,-.08) {$X,Y$ 可以在区域中连续移动\\无法把所有点对逐个列出来};
 \node[pnote,text width=6.2cm] at (11.55,-.08) {计算 $h(h-1)/2$ 个顶点对距离\\取最大值，并记下最远点对 $A,B$};
\end{tikzpicture}%
}

\newcommand{\CoverageSteps}{%
\begin{tikzpicture}[proc]
 \path[use as bounding box] (-.15,-.72) rectangle (15.05,4.18);
 \node[ptitle] at (2.15,3.93) {(a) 半径 $D/2$ 锁定圆心};
 \node[ptitle] at (7.45,3.93) {(b) 再逐个检查顶点};
 \node[ptitle] at (12.75,3.93) {(c) 漏点时，改求包围圆};
 \begin{scope}[shift={(2.15,1.92)}]
  \filldraw[procblue,fill=procblue!10] (-1,0) circle (1);
  \filldraw[procorange,fill=procorange!11] (1,0) circle (1);
  \draw[black!55] (-1,0)--(1,0);
  \fill (-1,0) circle (1.4pt) node[below] {$A$};
  \fill (1,0) circle (1.4pt) node[below] {$B$};
  \fill[procred] (0,0) circle (1.8pt) node[above] {$C_0$};
  \node[procblue,font=\scriptsize] at (-1,1.28) {$\|C-A\|\le D/2$};
  \node[procorange,font=\scriptsize] at (1,-1.28) {$\|C-B\|\le D/2$};
 \end{scope}
 \begin{scope}[shift={(7.45,1.92)},x=.85cm,y=.85cm]
  \filldraw[edge,fill=procteal!10] (-1.4,0)--(0,-1)--(1.4,0)--(0,1.7)--cycle;
  \draw[procred,dashed,thick] (0,0) circle (1.4);
  \fill[procred] (0,0) circle (1.5pt) node[below] {$C_0$};
  \fill[procred] (0,1.7) circle (2pt) node[above right] {漏点};
  \foreach \p in {(-1.4,0),(0,-1),(1.4,0)}{\fill[procteal] \p circle (1.5pt);}
  \node[below left] at (-1.4,0) {$A$};\node[below right] at (1.4,0) {$B$};
  \draw[procred,dotted] (0,1.4)--(.7,1.4) (0,1.7)--(.7,1.7);
  \draw[procred,<->] (.65,1.4)--(.65,1.7);
 \end{scope}
 \begin{scope}[shift={(12.75,1.92)},x=.85cm,y=.85cm]
  \filldraw[edge,fill=procteal!10] (-1.4,0)--(0,-1)--(1.4,0)--(0,1.7)--cycle;
  \draw[black!25,dashed] (0,0) circle (1.4);
  \draw[procblue,very thick] (0,.2735294118) circle (1.4264705882);
  \draw[procblue,->] (0,.2735294118)--(-1.4,0) node[midway,above] {$R^*$};
  \fill[procblue] (0,.2735294118) circle (1.6pt) node[right] {$C^*$};
  \foreach \p in {(-1.4,0),(0,-1),(1.4,0),(0,1.7)}{\fill[procteal] \p circle (1.7pt);}
 \end{scope}
 \ProcArrow{4.55}\ProcArrow{9.85}
 \node[pnote] at (2.15,-.23) {圆心必须同时落在两个圆盘内\\两圆盘的交集\textbf{只有中点}};
 \node[pnote] at (7.45,-.23) {以 $C_0$ 作圆，检查全部顶点\\有一个顶点在外，就不能覆盖};
 \node[pnote] at (12.75,-.23) {枚举两点圆与三点圆\\在覆盖全部顶点者中取最小};
\end{tikzpicture}%
}

% END PROCESS DRAWING DEFINITIONS

\begin{document}\thispagestyle{plain}
\begin{center}{\LARGE\bfseries 基于角域交会与最小包围圆的\par\vspace{0.25em}无源测向快速定位模型\par}\end{center}
% BEGIN ORIGINAL BODY WITH PROCESS FIGURES

\section*{摘要}
\addcontentsline{toc}{section}{摘要}

针对无线电干扰源的快速自动定位问题，本文建立角域交会定位模型，给出定位区域及其直径的求解算法，并对“以区域直径为直径的圆能否覆盖定位区域”这一判断作出回答。

\textbf{问题一}要求由若干检测点及其示向度计算干扰源的定位区域与直径，并判断直径圆的覆盖性。本文采用\textbf{角域交会定位模型}：示向度存在 $\pm1^\circ$ 的读数误差，一次测向只能确定干扰源位于一个张角 $2^\circ$ 的角域内；该角域可用两个线性不等式精确刻画，故多次测向的定位区域是若干半平面之交，为\textbf{凸多边形}。利用凸性，区域直径必在最远顶点对上取得，将区域上的极值问题化为有限次顶点比较，算法总复杂度 $O(n^3)$。对覆盖性，本文证明直径恰为 $D$ 的覆盖圆没有选择圆心的余地---圆心只能取最远点对的中点，由此得到 $O(h)$ 的充要判据；并构造了一个由两个误差角域直接求交得到的算例，其中该圆漏掉区域的一个顶点约 $0.307$ m，说明\textbf{以区域直径为直径的圆不一定能覆盖定位区域}。最后指出保证覆盖应采用\textbf{最小包围圆模型}，并以 $R^*\le20$ m 作为进入清除的判据。

\textbf{关键词：} 角域交会　半平面交　凸多边形直径　最小包围圆　示向度误差

\clearpage
\section*{一、问题背景与重述}
\addcontentsline{toc}{section}{一、问题背景与重述}

无线电干扰源的定位与清除是频谱管理中的基础任务。完整的定位流程通常包括初步监测、区域性排查与精准定位三个阶段，前两个阶段只能划定干扰源的大致分布范围，最终的精确定位仍需人员或设备抵近完成。为缩短排查时间、降低人员风险，本题考虑以机器狗代替技术人员，携带测向机与光学探测设备在目标区域内自动搜寻并清除干扰源。

目标区域是半径 $1800$ m 的圆形区域，圆心为坐标原点，正东为 $x$ 轴正方向，正北为 $y$ 轴正方向，单位为米。区域内分布 $10$ 至 $16$ 个干扰源，具体个数未知，每个干扰源占用一个互不相同的频道。

机器狗通过测向机获得的不是干扰源的位置，而是\textbf{示向度}：在检测点处测得的干扰源最大场强方向的方位角，自 $x$ 轴正方向逆时针计量，范围 $[0^\circ,360^\circ)$。受电磁环境与测向精度限制，示向度读数与真实方位之间存在误差，误差界限为 $[-1^\circ,1^\circ]$。也就是说，一次测向只能说明干扰源位于一个很窄的楔形范围内，要确定它的位置，必须从不同方向多次测向，逐步逼近---这正是\textbf{交会定位}的思路。

题目给出的问题一如下：

\begin{quote}
已知若干个检测点坐标及某干扰源关于这些检测点的示向度，请给出采用交会定位法计算多边形定位区域直径（即区域内任意两点之间距离的最大值）的算法。以定位区域直径为直径的圆能否覆盖此定位区域？
\end{quote}

本文即围绕这一问题展开。

\section*{二、问题一的分析}
\addcontentsline{toc}{section}{二、问题一的分析}

问题一包含两件事：一是算出定位区域及其直径，二是判断直径圆能否盖住这个区域。本问的难点不在数值计算，而在于找到一个合适的几何结构，使“区域内任意两点距离的最大值”这样一个涉及无穷多点的量能够有限步算出。

注意到每次测向给出的角域是一条射线张开得到的扇形，边界是两条直线，可以精确地写成两个线性不等式。因此，$n$ 次测向的定位区域就是 $2n$ 个\textbf{半平面的交}，是一个\textbf{凸多边形}。凸性带来了很大的便利：区域内任意一点都可以表示为顶点的凸组合，于是任意两点的距离不会超过最远顶点对的距离，而顶点之间的比较是有限次的。这样，求直径的问题就归结为求顶点、再枚举顶点对。

第二问的判断，关键在圆心。若一个圆的直径恰好等于区域直径 $D$，它必须同时覆盖最远的两点，由三角不等式可知圆心只能取这两点的中点，不能再作调整。圆心定下来后，覆盖与否只需逐一检验各顶点是否落在圆内。本文构造了一个完全符合题设误差界的算例，其中该圆没有盖住区域的一个顶点，说明第二问的答案是否定的；并指出此时应改用\textbf{最小包围圆}。

\Needspace{10\baselineskip}
\section*{三、模型假设}
\addcontentsline{toc}{section}{三、模型假设}

\begin{enumerate}
\item 测向误差有界，界为 $\pm1^\circ$；除误差界外，不对误差的分布形式作任何假设；
\item 检测点位置精确已知，问题在平面内讨论，忽略机器狗与干扰源的高度差异；
\item 各测向数据相容，即真实干扰源同时位于所有误差角域内；
\item 干扰源为点源，不考虑其发射功率与天线方向性对定位的影响。
\end{enumerate}

\Needspace{21\baselineskip}
\section*{四、符号说明}
\addcontentsline{toc}{section}{四、符号说明}

\begin{center}
\renewcommand{\arraystretch}{1.3}
\begin{tabularx}{0.98\linewidth}{>{\centering\arraybackslash}p{0.25\linewidth}>{\centering\arraybackslash}p{0.24\linewidth}>{\centering\arraybackslash}X}
\toprule
\textbf{符号} & \textbf{范围} & \textbf{说明} \\
\midrule
$n$ & $\mathbb{Z}^+$ & 测向次数（检测点个数） \\
$S_i=(x_i,y_i)$ & $\mathbb{R}^2$ & 第 $i$ 个检测点的坐标 \\
$\theta_i$ & $[0^\circ,360^\circ)$ & 第 $i$ 次测向的示向度 \\
$\varepsilon$ & $1^\circ$ & 测向误差界 \\
$P$ & --- & 定位区域 \\
$v_1,\dots,v_h$ & $\mathbb{R}^2$ & 定位区域的顶点 \\
$D$ & $(0,+\infty)$ & 定位区域直径 \\
$C^*,\ R^*$ & --- & 最小包围圆的圆心与半径 \\
\bottomrule
\end{tabularx}
\end{center}

\section*{五、模型的建立与求解}
\addcontentsline{toc}{section}{五、模型的建立与求解}

\subsection*{5.1 角域交会定位模型}
\addcontentsline{toc}{subsection}{5.1 角域交会定位模型}

设机器狗在 $n$ 个位置 $S_i=(x_i,y_i)$ 处对该干扰源测得示向度 $\theta_i$。$\theta_i$ 是自正东逆时针量到“检测点指向干扰源”方向的角度，但它并不是干扰源的真实方位：真实方位落在 $\theta_i$ 两侧各 $1^\circ$ 内。记 $\varphi_i=\pi\theta_i/180$，$\varepsilon=1^\circ=\pi/180$，则该角域的两条边界射线方向为

\[
u_i^-=\begin{pmatrix}\cos(\varphi_i-\varepsilon)\\ \sin(\varphi_i-\varepsilon)\end{pmatrix},\qquad
u_i^+=\begin{pmatrix}\cos(\varphi_i+\varepsilon)\\ \sin(\varphi_i+\varepsilon)\end{pmatrix},
\]

分别对应真实方位取到下限与上限两个极端。对二维向量记叉积 $a\times b=a_xb_y-a_yb_x$：若 $a\times b>0$，则 $b$ 在 $a$ 的逆时针一侧；若小于零，则在顺时针一侧；等于零则两向量共线。据此，点 $X=(x,y)$ 落在该角域内等价于

\[
u_i^-\times(X-S_i)\ge 0,\qquad u_i^+\times(X-S_i)\le 0. \tag{1}
\]

式 (1) 中第一式要求 $X-S_i$ 在 $u_i^-$ 的逆时针一侧，即位于下边界的上方；第二式要求它在 $u_i^+$ 的顺时针一侧，即位于上边界的下方。两式合起来，正是把 $X$ 限制在两条射线之间。又因角域张角仅 $2^\circ$，两条射线方向相反的那个区域（检测点背后）不满足这两式，不必另行排除。


% BEGIN PROCESS FIGURE
\begin{center}\begin{minipage}{\linewidth}
\noindent 图~\ref{fig:direction-process}~将示向度、误差角域与两条约束对应起来。
\par\medskip\centering
\DirectionToConstraints
\captionof{figure}{从示向度到半平面约束的转换。误差角作放大示意；实际计算仍取 $\varepsilon=1^\circ$。}
\label{fig:direction-process}
\end{minipage}\end{center}
% END PROCESS FIGURE

将式 (1) 展开，一次测向对应两个线性不等式：

\[
\begin{aligned}
\sin(\varphi_i-\varepsilon)(x-x_i)-\cos(\varphi_i-\varepsilon)(y-y_i)&\le 0,\\
-\sin(\varphi_i+\varepsilon)(x-x_i)+\cos(\varphi_i+\varepsilon)(y-y_i)&\le 0.
\end{aligned}
\tag{2}
\]

这两个式子都是 $x$、$y$ 的一次不等式，形如 $ax+by\le c$。满足 $ax+by\le c$ 的点集是直线 $ax+by=c$ 一侧的半个平面，称为半平面，它是一个凸集。式 (2) 的两个不等式分别限制了角域的两条边界，它们的交集就是该角域。

$n$ 次测向后，干扰源必须同时落在每个角域内，也就是要同时满足全部 $2n$ 个不等式。将它们统一写成 $a_jx+b_jy\le c_j$，得到定位区域

\[
P=\bigcap_{j=1}^{2n}\left\{(x,y):a_jx+b_jy\le c_j\right\}. \tag{3}
\]

由式 (3)，$P$ 是 $2n$ 个半平面的交。半平面是凸集，凸集之交仍是凸集，故 $P$ 为凸集；在非空且有界时，$P$ 是\textbf{凸多边形}，退化情形为线段或点。题目图 2 中的红色四边形，正是 $n=2$ 时两个角域求交的结果；多个检测点的情形依式 (3) 继续求交即可，每增加一次有效测向，约束增加两条，定位区域随之收缩。


% BEGIN PROCESS FIGURE
\begin{center}\begin{minipage}{\linewidth}
\noindent 图~\ref{fig:intersection-process}~按测向信息增加的顺序，展示公共区域怎样形成并继续收缩。
\par\medskip\centering
\ProgressiveIntersection
\captionof{figure}{逐次求交的过程：每次只保留已有区域中同时满足新约束的部分。图中角度和坐标为过程示意。}
\label{fig:intersection-process}
\end{minipage}\end{center}
% END PROCESS FIGURE

\subsection*{5.2 定位区域及其直径的求解}
\addcontentsline{toc}{subsection}{5.2 定位区域及其直径的求解}

求解分三步：先判明区域是否存在且有界，再求出全部顶点，最后在顶点对中取距离最大者。

\textbf{第一步，判别有界性。} 并非任意一组测向都能定出有限的区域，两种情形需要先排除。其一，若各次测向来自机器狗几乎沿一条直线走过的位置，各示向度方向接近平行，角域之交会沿这条方向无限延伸，区域无界。其二，若某次读数受干扰而偏差较大，各角域可能互不相容，区域为空。为此用线性规划对约束组作两个判断：先检验是否可行，不可行则说明数据与误差界矛盾，应剔除异常的测向数据；再对可行情形求 $x$、$-x$、$y$、$-y$ 在约束下的最大值，四者均有限当且仅当 $P$ 有界，若某一方向无界，则需改变探测位置，使各次测向的方向张开一些再重新定位。

有界性判别不能省略。无界区域同样可能有若干有限顶点，例如上述楔形，其尖端是有限顶点，但区域内沿楔形方向可以走到任意远；若跳过判别直接取顶点对的最大距离，会把本应为无穷的直径误报成有限值。


% BEGIN PROCESS FIGURE
\begin{center}\begin{minipage}{\linewidth}
\noindent 图~\ref{fig:feasibility-process}~说明为什么必须先分类，再进入有限直径的计算。
\par\medskip\centering
\FeasibilityCases
\captionof{figure}{可行性与有界性决定下一步操作。空集与无界集都不能直接套用有限顶点对的直径公式。}
\label{fig:feasibility-process}
\end{minipage}\end{center}
% END PROCESS FIGURE

\textbf{第二步，求顶点。} 凸多边形的每个顶点都是某两条边界直线的交点。对直线 $a_jx+b_jy=c_j$ 与 $a_kx+b_ky=c_k$，记 $\Delta_{jk}=a_jb_k-a_kb_j$，当 $\Delta_{jk}\ne0$ 时交点为

\[
x_{jk}=\frac{c_jb_k-b_jc_k}{\Delta_{jk}},\qquad
y_{jk}=\frac{a_jc_k-c_ja_k}{\Delta_{jk}}. \tag{4}
\]

两直线平行时无交点或有无穷多交点，都不产生顶点，故只需处理 $\Delta_{jk}\ne0$ 的情形。求出的交点未必在区域内，需代回全部 $2n$ 个约束逐一检验，保留满足 $a_\ell x_{jk}+b_\ell y_{jk}\le c_\ell$（$\ell=1,\dots,2n$）者；这些点还可能有重复（多条边界直线共点），去重后即得顶点集 $\{v_1,\dots,v_h\}$。


% BEGIN PROCESS FIGURE
\begin{center}\begin{minipage}{\linewidth}
\noindent 图~\ref{fig:vertex-process}~把交点处理展开为“枚举、约束检验、去重”三个动作。
\par\medskip\centering
\VertexExtraction
\captionof{figure}{顶点提取的三个步骤。图中用一个示意多边形展示绘图窗口内的候选交点；橙色空心点为候选，红色叉号为不合格点，绿色实心点为保留点。}
\label{fig:vertex-process}
\end{minipage}\end{center}
% END PROCESS FIGURE

\textbf{第三步，求直径。} 按题目定义，定位区域直径是区域内任意两点间距离的最大值，$D=\max_{X,Y\in P}\|X-Y\|_2$。由于 $P$ 是凸多边形，该最大值必在顶点对上取得：

\[
D=\max_{1\le i<j\le h}\|v_i-v_j\|_2 . \tag{5}
\]

下证式 (5)。凸多边形是其全部顶点的凸包，因此 $P$ 中任意一点都可写为顶点的凸组合，即存在权重 $\lambda_i\ge0$、$\sum_i\lambda_i=1$，使 $X=\sum_i\lambda_iv_i$；同理 $Y=\sum_j\mu_jv_j$，$\mu_j\ge0$、$\sum_j\mu_j=1$。两式相减得

\[
X-Y=\sum_{i,j}\lambda_i\mu_j(v_i-v_j),
\qquad
\sum_{i,j}\lambda_i\mu_j=\left(\sum_i\lambda_i\right)\left(\sum_j\mu_j\right)=1 .
\]

由三角不等式，

\[
\|X-Y\|\le\sum_{i,j}\lambda_i\mu_j\|v_i-v_j\|\le\max_{i,j}\|v_i-v_j\|\cdot1=\max_{i,j}\|v_i-v_j\|,
\]

即区域内任意两点的距离都不超过最远顶点对的距离。另一方面，顶点本身属于 $P$，故最远顶点对的距离是区域内可以取到的点对距离。两者结合，式 (5) 成立。这一结论把区域内无穷多组点对的比较，归结为顶点之间的有限次比较。


% BEGIN PROCESS FIGURE
\begin{center}\begin{minipage}{\linewidth}
\noindent 图~\ref{fig:diameter-process}~将式 (5) 的结论转化为可以逐项执行的距离比较。
\par\medskip\centering
\DiameterReduction
\captionof{figure}{由连续点对转为有限顶点对。灰线表示待比较的顶点连线，红线表示最远顶点对，其长度即区域直径。}
\label{fig:diameter-process}
\end{minipage}\end{center}
% END PROCESS FIGURE

综上，问题一的求解流程为：

\begin{tcolorbox}[colback=black!2,colframe=black!22,boxrule=0.4pt,arc=0pt,left=9pt,right=9pt,top=7pt,bottom=7pt]
\small\setlength{\parskip}{2pt}
\noindent\hangindent=1.5em\hangafter=1 输入：检测点 $S_i$、示向度 $\theta_i$，$i=1,\dots,n$\par
\noindent\hangindent=1.5em\hangafter=1 1  按式 (2) 将每次测向化为两个半平面；\par
\noindent\hangindent=1.5em\hangafter=1 2  线性规划检验可行性与有界性，异常情形分别报告；\par
\noindent\hangindent=1.5em\hangafter=1 3  按式 (4) 求每对边界直线的交点；\par
\noindent\hangindent=1.5em\hangafter=1 4  保留满足全部约束的交点并去重，得顶点集；\par
\noindent\hangindent=1.5em\hangafter=1 5  枚举所有顶点对，取距离最大者 $(A,B)$，$D=|AB|$。\par
\end{tcolorbox}

边界直线共 $2n$ 条，需要检验的交点约 $\binom{2n}{2}=O(n^2)$ 个，每个交点检验约束需 $O(n)$，顶点数不超过 $2n$，最后比较顶点对又是 $O(n^2)$，故总计算量 $O(n^3)$。机器狗每取得一次新的示向度就可用该算法更新一次定位区域，这一计算量在实时策略中不构成负担。

\subsection*{5.3 直径圆的覆盖性分析}
\addcontentsline{toc}{subsection}{5.3 直径圆的覆盖性分析}

机器狗携带的光学探测仪只能在距干扰源不超过 $20$ m 时精确定位并清除，因此“以定位区域直径为直径的圆能否覆盖该区域”，直接决定了能否将整片候选区域一次性纳入光学作用范围。

若半径为 $D/2$、圆心为 $C$ 的圆覆盖了 $P$，则 $A$、$B$ 两点都在圆内，即 $\|A-C\|\le D/2$、$\|C-B\|\le D/2$。代入三角不等式

\[
D=\|A-B\|\le\|A-C\|+\|C-B\|\le\frac D2+\frac D2=D .
\]

不等式的两端相等，故中间每一步都必须取等号：既要有 $\|A-C\|=\|C-B\|=D/2$，又要求三角不等式 $D\le\|A-C\|+\|C-B\|$ 取等，即 $C$ 在线段 $AB$ 上。两者合起来，$C$ 只能是 $AB$ 的中点 $C_0=(A+B)/2$。也就是说，直径恰为 $D$ 的圆，其圆心被 $A$、$B$ 唯一确定。又因为圆是凸集，它包含全部顶点就包含以这些顶点为顶点的整个凸多边形，所以覆盖 $P$ 等价于覆盖所有顶点，判据为

\[
\max_{1\le i\le h}\|v_i-C_0\|\le\frac D2 . \tag{6}
\]


% BEGIN PROCESS FIGURE
\begin{center}\begin{minipage}{\linewidth}
\noindent 图~\ref{fig:coverage-process}~将覆盖判据与失败后的处理连在一起；最小包围圆的求法在第 5.4 节展开。
\par\medskip\centering
\CoverageSteps
\captionof{figure}{覆盖判断的操作顺序：先锁定圆心，再检查每个顶点；若漏点，则改求最小包围圆。几何形状为定性示意，不采用表 1 的数值比例。}
\label{fig:coverage-process}
\end{minipage}\end{center}
% END PROCESS FIGURE

式 (6) 是否总成立？我们用下面的算例说明它可以不成立。取两个检测点

\[
S_1=(-500,\,0),\quad\theta_1=44.5^\circ;\qquad
S_2=(500,\,0),\quad\theta_2=135.5^\circ,
\]

误差界仍为 $\pm1^\circ$，两次测得的真实方位分别落在 $\theta_1$、$\theta_2$ 两侧各 $1^\circ$ 内。记 $t=\tan43.5^\circ$、$u=\tan45.5^\circ$，两个角域分别是 $t(x+500)\le y\le u(x+500)$ 与 $t(500-x)\le y\le u(500-x)$，按式 (3) 求交得一四边形，结果见表 1。

\Needspace{19\baselineskip}
\begin{center}
\textbf{表 1 算例的顶点与覆盖性检验（单位：m）}
\end{center}

\begin{center}
\renewcommand{\arraystretch}{1.3}
\begin{tabularx}{0.98\linewidth}{>{\raggedright\arraybackslash}p{0.52\linewidth}X}
\toprule
\textbf{量} & \textbf{数值} \\
\midrule
下顶点 $V_1$ & $(0,\ 474.482283)$ \\
右顶点 $V_2$ & $(17.452406,\ 491.043999)$ \\
上顶点 $V_3$ & $(0,\ 508.803696)$ \\
左顶点 $V_4$ & $(-17.452406,\ 491.043999)$ \\
直径 $D=\|V_2-V_4\|$ & $34.904813$ \\
圆心 $C_0=(V_2+V_4)/2$ & $(0,\ 491.043999)$ \\
半径 $D/2$ & $17.452406$ \\
上顶点到圆心的距离 $\|V_3-C_0\|$ & $\boldsymbol{17.759698}$ \\
\bottomrule
\end{tabularx}
\end{center}

枚举全部顶点对可知，最远的一对是左右两个顶点 $V_2$、$V_4$，距离 $D\approx34.905$ m，故圆心只能取 $C_0$、半径只能取 $D/2$。但上顶点到圆心的距离为

\[
\|V_3-C_0\|=508.803696-491.043999=17.759698\ \text{m}>17.452406\ \text{m},
\]

超出半径约 $0.307$ m，落在圆外，式 (6) 不成立。

从图形上看，原因在于这个四边形的两条对角线长度接近（$34.905$ m 与 $34.321$ m），但水平对角线并不平分竖直对角线：交点以上高 $17.760$ m，以下只有 $16.562$ m。圆心被限定在水平对角线的中点，位置偏低，故上顶点落在圆外。

该算例由两个满足题设误差界的角域直接求交得到，并非人为拼凑的图形；两次测向方向相差约 $91^\circ$，检测点相距 $1000$ m，与目标区域内的实际探测尺度一致；四边形顶点到原点最远约 $508.8$ m，整体位于半径 $1800$ m 的目标区域之内，故即便引入“干扰源位于目标区域内”这一先验，结论也不改变。于是问题一的第二问答案为否：

\textbf{以定位区域直径为直径的圆，不一定能覆盖该定位区域；即使允许自由选择圆心，结论也是如此。}

\subsection*{5.4 保证覆盖时的修正：最小包围圆模型}
\addcontentsline{toc}{subsection}{5.4 保证覆盖时的修正：最小包围圆模型}

上面的算例说明，“区域直径 $D$”与“盖住区域所需的最小圆半径”是两个不同的量。后者由最小包围圆模型给出：

\[
\min_{C,\ R}\ R
\qquad\text{s.t.}\quad\|v_i-C\|_2\le R,\quad i=1,\dots,h . \tag{7}
\]

其中，$C$ 是待求的圆心，$R$ 是待求的半径。由圆的凸性，只要全部顶点落在圆内，区域内的其余点也落在圆内，故只需对顶点施加约束。

对有限个点，最小包围圆有一条现成的性质：它或者以某两个顶点为直径的两端，或者经过某三个不共线的顶点。这是因为，若半径还可以继续减小，就说明圆没有被顶点“卡住”，沿某个方向平移圆心即可缩小半径；只有当圆被两个点（且圆心在两点连线上）或三个点同时约束时，半径才不能再减小。因此顶点不多时，枚举全部两点圆和三点圆，挑出能覆盖所有顶点且半径最小的一个，即为所求。

对表 1 的算例求解模型 (7)，得

\[
C^*\approx(0,\ 491.348632),\qquad
R^*\approx17.455065\ \text{m}>\frac D2\approx17.452406\ \text{m},
\]

即便把圆心放到最合适的位置，所需半径仍大于 $D/2$，与前面的判断一致。

进一步地，对一般的平面有界凸区域，两个量之间的差距有普适的界：

\[
\frac D2\le R^*\le\frac{D}{\sqrt3}. \tag{8}
\]

左端由最远点对给出。右端的来源是：若最小包围圆由三个顶点确定，则这三点构成锐角三角形，其最大内角 $\alpha$ 不小于 $60^\circ$，而由正弦定理，半径与边长的关系为 $R^*=\dfrac{\text{对边}}{2\sin\alpha}$，代入 $\alpha\ge60^\circ$ 即得 $R^*\le\dfrac{D}{\sqrt3}$；若由两点直径确定则更有 $R^*=D/2$。等边三角形恰好使上界取等，说明式 (8) 不能再改进，也就是说，直接把覆盖圆半径取成 $D/2$，最多可能低估约 $15.5\%$。

这个区分与后续清除策略直接相关。机器狗在某位置能否清除干扰源，取决于该位置与干扰源的距离是否不超过 $20$ m；若要保证当前位置覆盖干扰源的\textbf{全部}候选位置，应当检验 $R^*\le20$ m，而 $D\le40$ m 只是必要条件。问题三、四中，机器狗正是依据已获得的示向度不断收缩定位区域、减小 $D$ 与 $R^*$，待 $R^*$ 落入 $20$ m 后再发送清除指令；若只检查 $D$ 而忽略区域的形状，清除指令就可能落空。

\subsection*{5.5 模型说明}
\addcontentsline{toc}{subsection}{5.5 模型说明}

以上模型针对的是测向角域的交会区域。若进一步利用“干扰源必位于半径 $1800$ m 圆域内”这一先验，则更紧的候选区域为该圆域与 $P$ 的交，可能含有圆弧边界，不再必为多边形；此时仅枚举直线交点的顶点算法不再完备，需补充直线与圆交点、圆弧段的处理，不能以内接正方形等替代圆域后声称得到了精确直径。本文算例整体位于该圆域内部，故该先验不影响任何结论。

此外，实际探测中检测点还受有效接收半径（$1000$ m 至 $1500$ m）的限制：只有检测点到干扰源的距离不超过其有效接收半径时，机器狗才能获得该频道的示向度，因此可用于定位的检测点必然落在干扰源周围的一个环带内，测向方向也不会过于平行。这是实际数据的有利条件，但本文算法不依赖它，对任意给定的一组检测点与示向度均成立。

至于问题标题所问的“圆能否覆盖”，本文的回答是否定的，且原因不在于圆的大小不够，而在于\textbf{圆心被最远点对唯一确定}---即便换一个圆心、把半径放到最优，所需的圆仍可能比直径圆更大。

\Needspace{14\baselineskip}
\section*{六、问题一小结}
\addcontentsline{toc}{section}{六、问题一小结}

\begin{enumerate}
\item 每次测向对应一个张角 $2^\circ$ 的误差角域，恰为两个半平面约束；$n$ 次测向的定位区域是 $2n$ 个半平面之交，非空有界时为凸多边形；
\item 先判别有界性，再枚举边界直线交点得到顶点，区域直径等于最远顶点对的距离，算法总复杂度 $O(n^3)$；
\item 直径为 $D$ 的覆盖圆，圆心只能取最远点对的中点；构造的算例中该圆漏掉上顶点约 $0.307$ m，最小包围圆半径 $R^*\approx17.455$ m 亦大于 $D/2$。故直径圆不一定能覆盖定位区域，需要保证覆盖时应改用最小包围圆，并以 $R^*\le20$ m 作为光学精确定位可达的判据。
\end{enumerate}
% END ORIGINAL BODY
\end{document}
